\input amstex
\documentstyle{amsppt}
\magnification 1200
\hsize 6.5 true in
\vsize 8.7 true in
\define\eps{\epsilon}
\define\ooo{{\Cal O}}
\define\ovo{\overline{\Cal O}}
\define\off{\overline{F}}
\define\oee{{\Cal O}_E}
\define\VV{\Cal V}
\define\HH{\Cal H}
\define\slf{G'(F)}
\baselineskip 16pt
\document
\def\leaderfill{\leaders\hbox to 1em{\hss.\hss}\hfill}
\def\eleaderfill{\hfill \hbox to 1em{\hss.\hss}\hfill}
\centerline{{\bf On the transfer of Shalika germs for $G_2$}}
\vskip.1in
\centerline{by David Joyner}
\vskip.1in
\centerline{summer, 1994}
\vskip .7in
{\bf \S 0 Introduction}
\vskip .1in
{\bf Motivation and rough outline}
\vskip .2in
To apply the trace formula to the transfer of automorphic representations
of two reductive groups $G$, $H$ over a number field $K$, one must,
among other things, have a comparison of local orbital integrals, for each 
local completion $F$ of $K$. The comparison in the real and complex case 
was solved by Shelstad \cite{S1}. 
In the $p$-adic case, such a local comparison - the smooth matching conjecture
stated below - 
boils down, thanks to the work of R. Langlands and D. Shelstad
(\cite{LS1}, \cite{LS2}), to a certain matching of Shalika germs
conjecture, also stated below. 
Shalika germs are mysteriously behaving functions
associated to a semisimple element $s\in G(F)$ and a unipotent
conjugacy class ${\Cal O}$ of $M=Cent_{G(F)}(s)$, defined in a punctured 
compact-open neighborhood of $s$ in the regular set \cite{Sh}. They play a 
role in the analysis of $p$-adic orbital integrals on reductive groups
roughly analogous to the role that the power functions play in the Taylor 
series of a function on the real line. By a well-known reduction known
as Harish-Chandra descent, it usually suffices to assume that $s$ is central
and $M=G$. 
Langlands \cite{L} introduced a geometric approach to the computation,
or at least comparison, of Shalika germs. 
The comparison of Shalika germs associated to the
$0$-regular unipotent orbit was solved by Shelstad \cite{S2},
using results from \cite{LS1}, \cite{LS2}.
Langland's technique was further developed by
T. Hales in his thesis \cite{H1}. Hales concentrated on the comparison of 
$1$-regular germs, matching the stable subregular germs in general and
matching the $\kappa$-stable subregular germs for even rank unitary
groups. He also worked out a number of the computations necessary
to deal with the rank two case. This technique of Langlands goes by
the name Igusa theory since it was Igusa \cite{I} who developed a crucial
asymptotic formula for an integral over fibers of a morphism
$Y_C\rightarrow C$, where $C$ is a curve and $Y_C$ is a non-singular
variety. This technique also relies on a formula for the coefficients
of this expansion, involving principal value integrals over $p$-adic 
analytic manifolds, worked out in \cite{I}, \cite{L} (see also
\cite{LS0}). In \cite{L}, \cite{LS3}, Langlands and Shelstad
matched all the germs for groups of absolute type $A_2$
(see also \cite{H3}, where a descent argument is used instead of
the direct comparison in \cite{LS3}).
In \cite{H2}, Hales worked out the comparison of germs for the forms of 
$B_2=C_2$ using Igusa theory and a reduction/descent argument. 
In the present article, we verify the matching of Shalika germs conjecture
mentioned above for the $1$-, $2$- and $3$-regular germs for the 
exceptional group $G_2$,
which, together with the result of \cite{S2}, proves the smooth 
matching theorem for $G_2$. In other words, our main result
is the proof of the following
\proclaim{Smooth Matching Conjecture for $G_2$}
Let $T_G$ denote an elliptic Cartan of $G=G_2$ 
which has an $F$-embedding into $H$ and let $\kappa$ denote
a character of $H^1(Gal(\off/F),T_G)$ associated to a
cuspidal endoscopic group $H=H(T_G,\kappa)$
\cite{L} (thus $H$ is either $PGL(3)$ or $(SL(2)\times SL(2))/\{\pm 1\}$).
Let $T_H$ denote the image of an $F$-embedding of $T_G$ into $H$.
For each $f\in C_c^\infty(G(F))$ there is an $f^H\in C_c^\infty(H(F))$
for which
$$
\Delta_{G/H}^{T_G,\kappa}(\gamma_H,\gamma_\lambda)
\Phi^{T_G,\kappa}(\gamma_G,f)
=\Phi^{T_H,st}(\gamma_H,f^H),
$$
whenever $\gamma_G$ is the image of a regular $\gamma_H\in T_H(F)$.
(The notation is explained further below.) 
\endproclaim
\vskip .1in
Before sketching the ideas in the paper, we describe what is new
or different from \cite{H2}, \cite{H3}, which this paper
is comparable to. The notation will be explained later.
\par
In \S 1, we carefully recall
in detail the construction of the Igusa variety and the
determination of local coordinates \cite{L}, \cite{H1}, \cite{H2}.

Section 2 gives the calculations for determining the action of the Galois
group on the local coordinates on the
Igusa variety for $G_2$ as in \cite{H2}. 

Section 3 describes some new results: for any reductive group $G$, 
we produce specific patches generalizing those used in \cite{H2} 
and show for $G_2$
that they cover the base Igusa variety $Y''$ (described below) at least up to 
codimension two. These specific patches are needed for doing certain 
local calculations. 
It is likely that these patches are regular under more
general conditions, which would imply $Y''$ is regular in codimension one 
(conjectured in \cite{H1}), but this is only proven in our
special case. Anyway, by another method
we show that, as a consequence of the theory of toroidal embeddings \cite{K},
$Y''$ is normal and Cohen-Macauley (see Lemma 3.7 below).
This implies Hales' conjecture that 
$Y''$ is regular in codimension one in general. 

In section 4 we compute the Igusa constants of all
the base divisors on $Y''$ by an easy direct calculation
and give an ``endoscopic interpretation'' of some of the
calculations.
We also compute the Igusa forms using local coordinates of $Y_1$.

In section 5, we first record some basic facts about
the Igusa expansion over $Y_C$. Then we
prove the geometric ``no spurious divisors lemma'' for $G_2$
which is basically an inequality $e(\theta,\beta)\leq 1$,
where $e(\theta,\beta)$ is as in \cite{L}, which allows one to work
on $Y''$ rather than its desingularization $Y_C$.
Finally, we recall the statement of Hales' important ``reduction lemma''.

In sections 6, 7, and 8
we (a) compute coordinates
on the $1$-, $2$- and $3$-regular unipotent orbits (the latter two
are not Richardson which is something new, though they are what I call
``endoscopically Richardson''), (b) compute
coordinates on the fiber $\phi^{-1}(\gamma_\lambda)^\Gamma$
(explained below), and (c) compute the Igusa form
on the fiber and its explicit dependence on a tangent vector $X$
of $C$. Not much in these 
sections is substantially different from \cite{H2}, \cite{H3}. One
new idea is that, roughly speaking, the subgroup generated by the
one-parameter root-groups associated to roots ``outside $H$'' 
may be used to determine local coordinates for the
``endoscopically Richardson'' orbit $\ooo=Ind_H^G(1)$.

The last section presents a conjecture which our method of proof
seems to support. It seems to be the ``ultimate'' theorem which one
might prove by this method without introducing any new ideas.

For the most part, the appendices collect explicit 
group-theoretic computations 
needed to justify a number of claims made in the main body of the paper. 
For example, we verify that the cuspidal endoscopic groups of $G_2$ are
$PGL(3)$ and $(SL(2)\times SL(2))/\{\pm 1\}$. 
\par
Turning now to a brief discussion of the constructions we use,
here is a rough sketch of Langlands' idea in its simplest form: 
Let $\gamma\in G(F)$ be an elliptic regular element and let
$T_G=Cent_G(\gamma)^o$. The goal is to compute,
or at least find some manageable expression for, the stable orbital
integral $\int_{(T_G\setminus G)(F)}f(g\gamma g^{-1})dg$. 
(More generally, one can and does consider $\kappa$-stable orbital 
integrals.) The variety $T_G\setminus G$
may be identified, over the algebraic closure
$\overline{F}$, with the variety of regular stars: let
$\Omega$ denote the absolute Weyl group of $T_G$, 
$B_G$ a Borel subgroup of $G$
containing $T_G$, and $S_0$ denote the $G(\off)$-orbit of the
$|\Omega|$-tuple $(B_G^\omega)_{\omega\in \Omega}$, where
$G$ acts by conjugation on each coordinate. This orbit is called the
variety of regular stars and since 
$\cap_{\omega\in \Omega}N_G(B_G^\omega)=T_G$,
we have $S_0\cong T_G\setminus G$ over $\off$. Let
${\Cal B}\cong B_G\setminus G$
denote the Borel variety. Although ${\Cal B}$ is complete,
$S_0$ is not but
conjugation induces a principal $N_G$-bundle structure on
$S_0\rightarrow {\Cal B}$, where $N_G$ denotes 
the unipotent radical of $B_G$, 
so $S_0$ may be roughly regarded as 
the base space ${\Cal B}$ with fibers at each point isomorphic to
$N_G$. The Zariski closure of
$S_0$ in ${\Cal B}^{|\Omega|}$ is denoted $S$ and is called the variety of 
stars. It is projective, complete, but contains singularities.
It turns out that $S_0$ is Zariski open in $S$. 
Roughly speaking, the idea is to regard the above stable orbital integral
as an integral over $S_0(F)\subset S(\off)$. Let $T_G'$
denote the subvariety of regular elements of $T_G$. We 
restrict $\gamma$ to a regular curve $C$ in $T'_G\cap \{1\}$
in a neighborhood of the identity and then construct a smooth variety $Y_C$
with a fibering $\phi:Y_C\rightarrow C$, ``the''
smooth Igusa variety, whose fibers
$\phi^{-1}(\gamma)$, $\gamma\in C,\ \gamma\not= 1$,
are isomorphic to $T_G\setminus G$. In fact, the above stable orbital integral
may be regarded now as a fiber integral parameterized by a 
local coordinate $\lambda$ in a neighborhood of the singular point $1\in C$. 
Taking the limit as $\lambda$ tends to this singular point,
one may then use techniques of Igusa \cite{I}, as developed
by Langlands \cite{L}, to express the stable orbital integral
as a sum involving principal value integrals over certain divisors
contained in $\phi^{-1}(1) \subset Y_C$. By the homogeneity
of the Shalika germs and the
homogeneity of the coefficients in the Igusa expansion, only certain
divisors can contribute to a given Shalika germ.
\par
Here is a rough sketch of Igusa theory \cite{L}, \cite{H1}: 
Assume $G$ is $F$-quasisplit. Let $C$ denote a $1$-parameter subgroup of $G$
in $T_G$, defined over $F$, containing the identity
and no other singular point and whose tangent over $\off$ is not
contained in any root hyperplane. Let $C'=C-\{1\}$. We sketch the 
construction of a non-singular algebraic variety $Y_C$, defined over $F$,
for which
$$
\CD
C'\times T_G\setminus G @> i >> Y_C @> \pi >> G\\
@. @VV \phi V @. \\
@. C, @. 
\endCD
$$
with the following properties:
\newline
(a) if $T_G$ and $C$ are defined over $F$ then so are
$Y_C$, $\phi$, $\pi$, and $i$;
\newline
(b) $\pi : Y_C\rightarrow G$ is proper and $G$-equivariant;
\newline
(c) $\phi : Y_C\rightarrow C$ is proper;
\newline
(d) $\pi(\phi^{-1}(1))=V$, the unipotent variety of $G$;
\newline
(e) if ${\Cal O}\subset G$ denotes a 
unipotent orbit and if $E=\pi^{-1}(u)$, 
for some $u\in \pi(Y_C)\cap {\Cal O}$,
then $\pi|_{E}:E\cap \pi^{-1}({\Cal O})\rightarrow {\Cal O}$
is faithfully flat 
\cite{DG, Ch II, Proposition 3.1} and, in particular, surjective; 
\newline
(f) $i(C'\times T_G\setminus G
)\subset Y_C$ is open;
\newline
(g) $\pi\circ i(\gamma,g)=\gamma^g$ and $\phi\circ i(\gamma,g)=\gamma$.
\newline
(h)  $Y_C$ is $2p+1$-dimensional, where $p=|\Sigma^+|$ 
denotes the number of positive roots.
\newline
A divisor $E\subset Y_C$ for which $\pi(E)=\ovo$ is called an
{\bf $\ooo$-divisor}. In general, although $E$ is smooth,
$\ovo$ need not even be normal \cite{KP, \S 0}.
\newline
Step 1: Let $X$ denote the Zariski closure of $T'\times T_G\setminus G$
in $G\times S$ and let $Y$ denote the Zariski closure of
$C'\times T_G\setminus G$
in $X$. Let $\chi :G\rightarrow T_G/\Omega$ denote the adjoint
quotient \cite{Slo}, where $\Omega$ denotes the Weyl group of $T_G$.
We will see below that $X$ is covered by patches 
$X(B_+ ,B_{-})\subset
B_{-}\times S\times N_+$ and $Y$ is covered by patches 
$Y(B_+ ,B_-)\times N_+\hookrightarrow \chi^{-1}(C) N_-\times S\times N_+$,
where $N_-,N_+$ denotes the unipotent radical of $B_-,B_+$, resp..
Here $B_+,B_-$ are opposite Borels for which
$T_+=B_+\cap B_-$ is defined over $F$. The completeness
of $S$ implies that the projections $Y\rightarrow G$
and $X\rightarrow G$ are proper.     
\newline
Step 2: We construct below a resolution $S_1$ of $S$ with a proper
morphism $S_1\rightarrow S$ and an injective morphism
$S_0\rightarrow S_1$. The image of $S_0$ in $S_1$ is denoted $S''$.
Both $S$ and $S_1$ are complete but $S''$ is not.
Let $X''$, $X_1$ denote the Zariski closure of $T'\times T_G\setminus G$
in $G\times S''$, $G\times S_1$, respectively. 
Let $Y''$, $Y_1$ denote the Zariski closure of $C'\times T_G\setminus G$
in $X''$, $X_1$, respectively. Again, we have proper maps
$Y_1\rightarrow G$ and $X_1\rightarrow G$. Though $Y''\rightarrow G$
may not be proper, it is especially useful because
of its relatively simple local coordinate relations. We can explicitly
desingularize $Y''$ for $G=G_2$ by blowing up four
times in succession. 
(The derivation of this is lengthy and not really needed for the
proof of our main result so is omitted.)
Recall that the blow down morphisms
are proper and defined over $F$. They are also $G$-equivariant
by the universal property of blowing up \cite{Har}.
\newline
Step 3: Resolve the singularities of $Y_1$ contained in
$Y_1-Y''$ by using Mumford's semi-stable reduction theorem \cite{K}.
This yields a variety $\tilde{Y_1}$ and a blow down
$\xi:\tilde{Y_1}\rightarrow Y_1$ such that all the singularities 
of $\tilde{Y_1}$ map down to $Y_1- Y''$ under $\xi$. Now
perform an embedded resolution of singularities on  
$Y''\hookrightarrow \tilde{Y_1}$ as in step 2. The resulting
resolution of $\tilde{Y_1}$
is denoted $Y_C$. The blow down map $Y_C\rightarrow Y_1$
is proper, so $\pi$ and $\phi$ are proper. 
\par
We remark that the properness of $\pi$ and
$\phi$ are needed to guarantee that the appropriate
principal value integrals make sense. Properness implies that 
the principal value integral of the
pull-back of the function $f_\kappa$ is well-defined
and that some of the hypotheses to Lemma 5.7 below hold.

\vskip .1in
{\bf Notation and detailed outline}
\vskip .1in
Let $\gamma_\lambda=exp(\lambda X)\in C'(F)\subset T_G(F)$ be the image of
$\gamma_{\lambda,H}\in T_H(F)\subset H(F)$ in the sense of \cite{LS2}.
After some general background, we will examine the $1$-, $2$-, 
and $3$-regular germs for $G_2$ using the formula (\cite{Sh}, \cite{LS2})
$$
\split
\Delta_{G/H}^{T_G,\kappa}(\gamma_{\lambda,H},\gamma_{\lambda})
\Phi^{T_G,\kappa}(\gamma_{\lambda},f)
&=\kappa(inv(\gamma^0_{G},\gamma_{\lambda}))
D_H(\gamma_{\lambda,H})^{-1}
\int_{\phi^{-1}(\gamma_{\lambda})^\Gamma}\pi^*(f)
\kappa(t_*){|\omega|\over |\phi^*(d\lambda)|}\\
&=\sum_{{\Cal O}}\mu_{\Cal O}(f)
\Gamma_{{\Cal O}}^{T_G,\kappa}(\gamma_{\lambda}),
\ \ \ f\in C_c^\infty(G),
\endsplit
\tag 0.1
$$
where $\Delta_{G/H}^{T_G,\kappa}$ is the local transfer factor given in a
neighborhood of the identity by 
$$
\Delta_{G/H}^{T_G,\kappa}(\gamma_{\lambda,H},\gamma_G)
=\kappa(inv(\gamma^0_G,\gamma_G))D_G(\gamma_G)/D_H(\gamma_H)
$$
as in \cite{LS2}, $\gamma_G^0$ being a fixed regular element of $T_G(F)$.
The sum over ${\Cal O}$ in (0.1) runs over all stable unipotent orbits
${\Cal O}^{st}$ in $G(F)$, ${\Cal O}\subset {\Cal O}^{st}$
denoting an arbitrarily chosen $G(F)$-orbit within ${\Cal O}^{st}$.
Also, in (0.1),
$\gamma_\lambda=\gamma_{\lambda,G}\in C\cap V_f\subset T_G$ 
has parameter $\lambda$,
so that $\gamma_0=1$, where $V_f\subset G_{reg}$
is a punctured neighborhood of $1$,
$t_*$ is a cocycle-valued function which will be defined following
0.6 below, and
$\phi^{-1}(\gamma_\lambda)^\Gamma$ denotes the fixed points of the fiber 
under the action of $\Gamma:=Gal(\overline{F}/F)$ on $Y_C$ determined
by $T_G$ as described in \S 2 below.
The measure
${|\omega|\over |\phi^*(d\lambda)|}$ incorporates the factor
$D_G(\gamma_{\lambda,G})$.
If $\gamma_G'\in G(F)$ is stably conjugate to $\gamma_G$ then
$\Delta_{G/H}^{T_G,\kappa}(\gamma_H,\gamma'_G)
=\kappa(inv(\gamma_G,\gamma'_G))
\Delta_{G/H}^{T_G,\kappa}(\gamma_H,\gamma_G)$.
We will normalize
the measures and use the same notation as in \cite{LS3}. In particular,
$$
\Gamma_{{\Cal O}}^{T,\kappa}(\gamma)=
\Delta_{G/H}^{T,\kappa}(\gamma_H,\gamma)\sum_\delta \kappa(\delta)
\Gamma_{\Cal O}(\gamma^\delta),
$$
if $\gamma$ is an image of some $\gamma_H$. Here $\delta$ runs over 
classes in $H^1(\Gamma,T_G)$ associated to the distinct stable 
conjugacy classes $\{\gamma^\delta\}$ of $\gamma$. These germs
$\Gamma_{\Cal O}(\gamma)$ are slightly different from those in \cite{Sh}
since the germs here do not include the discriminant factor
in their definition. By a result of Harish-Chandra, if
${\Cal O}$ is $r$-regular then $\Gamma_{\Cal O}$ satisfies
$$
\Gamma_{\Cal O}(exp(\lambda^2 X))=|\lambda|^{-dim\ {\Cal O}}
\Gamma_{\Cal O}(exp( X)),
$$
where for any $u\in {\Cal O}$, 
$dim\ {\Cal O}=dim\ G_u\setminus G=dim\ G -rk(G)-2r=2|\Sigma^+_G|-2r$ with
$G_u$ denoting the centralizer of $u$ in $G$, 
$r=dim\ {\Cal B}_u$, ${\Cal B}$ denoting the Borel variety,
${\Cal B}_u$ its subvariety fixed by the conjugation action of $u$
and $\Sigma^+_G$ denoting the positive roots with respect to some Borel.
Note that since
$$
D_G(exp(\lambda X))=\prod_{\sigma >0}(1-\sigma(exp(\lambda X))^{-1})
=\lambda^{|\Sigma^+_G|}\prod_{\sigma >0}\sigma(X)+\text{ higher order terms},
$$
the normalized germ $\Gamma^*_{\Cal O}(\gamma)
:=|D_G(\gamma)|\Gamma_{\Cal O}(\gamma)$ satisfies
$\Gamma^*_{\Cal O}(exp(\lambda^2 X))
=|\lambda|^{2r}\Gamma^*_{\Cal O}(exp( X))$.

The ultimate object of our desires is the following
\proclaim{Matching of Germs Conjecture \cite{LS3}} 
Let $H$ be an endoscopic group 
associated to $T,\kappa$. Given an $r$-regular
unipotent orbit ${\Cal O}$ in $G(F)$ there is a finite collection 
of $r$-regular unipotent classes ${\Cal O}_i$ in $H(F)$ and complex
numbers $a_i$ such that, for every Cartan $T_H$ of $H$ over $F$ and every
regular ray $\{\lambda X|\lambda\not=0\}$ in the Lie algebra ${\frak t}_H$,
we have
$$
\Gamma_{{\Cal O}}^{T,\kappa}(\gamma_\lambda)=
\sum_i a_i\Gamma_{{\Cal O}_i}^{T_H,st}(\gamma_\lambda ),
$$
for $\lambda$ sufficiently small.
\endproclaim
Our goal is to try to verify this for $G=G_2$ and $r=1,2,3$.
\par
The expression in (0.1) may also be written in
the form $1/D_H(\gamma_{\lambda,H})$
times
$$
F(\lambda)=\int_{\phi^{-1}(\lambda)(F)}f_\kappa(y)|\omega_\lambda|,
\tag 0.2
$$
where $f_\kappa$, as defined by (0.1) and (0.2),
satisfies the conditions of being an Igusa integrand
in the sense of \cite{L}. Igusa's expansion
\cite{L} is an asymptotic expression of the form
$$
F(\lambda) =\sum_{\theta,\beta}\theta(\lambda)|\lambda|^{\beta -1}
\sum_{r=1}^{e(\theta,\beta)}(-\log_q|\lambda|)^{r-1}F_r(\theta,\beta,f_\kappa),
\tag 0.3
$$
where $\theta$ is a unitary character of $F^{\times}$ and $\beta >0$
is a rational number for which ${\Cal E}(\theta,\beta)\not= \emptyset$.
Here ${\Cal E}(\theta,\beta)$ denotes the set of irreducible components
(called {\bf branches}) in $\phi^{-1}(1)\subset Y_C$ such that
$\theta^{a(E)}=\kappa_E$ and $\beta=b(E)/a(E)$, where $a(E)$, $b(E)$, 
$\kappa_E$, and $e(\theta,\beta)$ are as in \cite{L}. As T. Hales
has observed, since there are no logarithmic terms in the
Shalika germ expansion, we have 
\proclaim{Lemma 0.4} If $f_\kappa$ is as in (0.2) then
$F_r(\theta,\beta,f_\kappa)=0$
for all $r>1$. 
\endproclaim
By Howe, Harish-Chandra, and
Rogawski \cite{Ro}, the trivial germ is a non-zero constant $c_H$.
We describe how we will be investigating the other
terms in (0.1) associated to the 
$1$-regular, $2$-regular,
and $3$-regular orbits ${\Cal O}$.
These will match with the trivial orbit in the germ expansion on $H$,
$$
\Phi^{T_H,st}(\gamma_{\lambda,H},f^H)
=\int_{\phi_H^{-1}(\gamma_{\lambda,H})^\Gamma}\pi^*(f^H)
{|\omega|\over |\phi^*(d\lambda)|}
=\sum_{{\Cal O}}\mu_{\Cal O}(f^H)
\Gamma_{{\Cal O}}^{T_H,st}(\gamma_{\lambda,H}),
\tag 0.5
$$
where $H=(SL(2)\times SL(2))/\{\pm 1\}$ or $H=PGL(3)$, and where
$f^H\in C_c^\infty(H)$ exists and
$\Delta^{T_G,\kappa}_{G/H}$-matches with $f$.
These terms are the $\beta =2,3,4$ terms in the Igusa expansions
for $G_2$, $H_1=PGL(3)$ and $H_2=(SL(2)\times SL(2)))/\{\pm 1\}$.
It ultimately turns out that, in these expansions, we have
$e(\theta,\beta)\leq 1$ (this is proved in \S\S 4-5 below).
Therefore, formula 0.6 below applies to our situation.
\par
Putting together results of \cite{L}, \cite{H2, \S 2},
using the Shalika germ expansion, we
have the following 
\proclaim{Formula 0.6(a)} Assume that we have a Cartan $T_G$, a character 
$\kappa$ of $H^1(\Gamma,T_G)$, and
a $(\beta-1)$-regular
stable unipotent class ${\Cal O}^{st}$ for which there is only
one $F$-rational irreducible branch $E_\beta$ of $\{\lambda =0\}$ in
$Y_C$ with $\pi(E_\beta)=\overline{{\Cal O}}^{st}$ and 
$\beta(E_\beta)=\beta$. Then we have
$$
D_H(\gamma_{H,\lambda})
\sum_{{\Cal O}\subset {\Cal O}^{st}}
\mu_{{\Cal O}}(f)\Gamma_{{\Cal O}}^{T_G,\kappa}(\gamma_\lambda)
=\kappa(inv(\gamma_G^0,\gamma_\lambda))
|\lambda|^{\beta -1}\sum_{\theta}\theta(\lambda)
\int_{E_\beta(F)} \pi^*(f)\kappa_\theta(t_*)|\omega_{E_\beta}|,
$$
for $\gamma_\lambda\in C\cap V_f$. Here $\kappa_\theta(t_*)$
is an extension of $\kappa(t_*)/\theta(\lambda)$ to
$E_\beta$ and $\overline{{\Cal O}}^{st}$ denotes the $F$-Zariski closure
of the stable orbit and $t_*$ is a cocycle-valued function 
described more below, and the sum on the left runs over all
unipotent $G(F)$-orbits ${\Cal O}$ contained in the
stable orbit ${\Cal O}^{st}$.
\newline
0.6(b): Same assumptions. Then for each $\overline{\Cal O}^{st}$
there is a measure $\omega_{\overline{\Cal O}^{st}}$ and for each
$F$-rational fiber
$E_\beta(u)=E_\beta(F)\cap \pi^{-1}(u)$,
$u\in \overline{\Cal O}(\overline{F})$,
there is a measure
$\omega_{E_\beta(u)}$ such that
$$
\int_{E_\beta(F)} \pi^*(f)\kappa_\theta(t_*)|\omega_{E_\beta}|
=\int_{\overline{\Cal O}^{st}}f\int_{E_\beta(u)} \kappa_\theta(t_*)
|\omega_{\overline{\Cal O}^{st}}||\omega_{E_\beta(u)}|,
$$
for each $\theta$ occuring above.
\endproclaim

This will follow from the results of section 5 below.
By \cite{H1, ch 6, \S 2}, only the $\theta =1$ 
term on the left side of formula 0.6(a) contributes a non-zero
term to the sum.
It turns out that either $\kappa$ is 
quadratic (and $H=(SL(2)\times SL(2))/\{\pm 1\}$)
or $\kappa$ is a cubic character (and $H=PGL(3)$).

Now we turn to recalling some of the properties of $t_*$ \cite{L},
\cite{H1}. Recall $\Gamma=Gal(\overline{F}/F)$.
The cocycle-valued function $t_*$ is a map
$Y_C\rightarrow Z^1(\Gamma, T_G(\overline{F}))$ given by
$P\longmapsto (\sigma\longmapsto t_\sigma(P))$
which has the following properties:
\newline
(a) it is locally constant on the $p$-adic manifold $Y_C(F)$ \cite{L};
\newline
(b) it factors through $Z^1(Gal(E/F), T_G(E))$,
where $E$ denotes the Galois closure of the splitting
field of $T_G$;
\newline
(c) for fixed $\sigma\in \Gamma$, $t_\sigma$ is a regular on the variety
$Y_C$;
\newline
(d) for $y=i(\gamma_\lambda ,g)\in Y_C$, such that for all
$\sigma\in \Gamma$,  
$\sigma(g)g^{-1}$ belongs to $T_G(\overline{F})$, we have
$\kappa (t_*(y))=\kappa(\sigma(g)g^{-1})$, where
$\sigma(g)g^{-1}$ represents
the cocycle $\sigma\longmapsto \sigma(g)g^{-1}$. 
\par
After $Y_C$ has been constructed, the idea is to determine the sets
${\Cal E}(\theta, \beta)$ with $\beta =2,3,4$. First,
one determines the $\off$-branches in $\phi^{-1}(1)\subset Y_C(\off)$ which
map down to $Y''$ under $Y_C\rightarrow Y_1$ which have
$\beta =2,3,4$. These are called {\bf fundamental divisors}.
In fact, in the case of $G=G_2$, we can construct on $Y''$ itself
divisors occurring in $\phi^{-1}(1)\subset Y''$.
Such divisors will be called {\bf base divisors}.
Of course, a base divisor is fundamental. 
We sometimes identify a base divisor with its closure in $Y_C$.
Then one determines which of these fundamental divisors
has $F$-rational
coordinates with respect to the Galois action determined by $T_G$.
These branches will contribute to Igusa's
expansion. Next, one must consider the $\off$-branches in 
$\phi^{-1}(1)\subset Y_C(\off)$ which do not
map down to $Y''$ under $Y_C\rightarrow Y_1$. These are 
called {\bf spurious divisors} and are relatively complicated to deal with. 
It turns
out that the computations in \cite{H1, Ch III} show that there are 
no $F$-rational spurious divisors contributing to the subregular germ
for $G_2$. For other reasons, the spurious divisors can be avoided
altogether for the $2$-regular and $3$-regular germs.
In fact, we have the following result, proven in \S\S 4-5 below.
\proclaim{Lemma 0.8} Let $E$ be an $F$-rational 
base ${\Cal O}$-divisor of the 
variety $Y''$ and let $\tilde{E}$ denote its strict transform
in $Y_C$. The principal
value integrals $\int_{{E}(F)}\pi^*(f)\kappa_1(t_*) |\omega_{{E}}|$ and
$\int_{\tilde{E}(F)}\tilde{\pi}^*(f)\kappa_1(\tilde{t}_*) |\omega_{\tilde{E}}|$
exist. Here $\kappa_1$ is as in formula 0.6(a), $t_*$ is the
cocycle function on $Y''$, and $\tilde{t}_*$ is the cocycle function
on $Y_C$. Moreover $e(1,\beta_E)=1$ and
$F_1(1,\beta_E,f_\kappa)$ is given 
by $\kappa(inv(\gamma_G^0,\gamma_G))$ times
$$
\split
PV\int_{{E}(F)}\pi^*(f)\kappa_1(t_*) |\omega_{{E}}|
&=PV\int_{\tilde{E}(F)}\tilde{\pi}^*(f)
\kappa_1(\tilde{t}_*) |\omega_{\tilde{E}}|\\
&=\int_{\overline{\Cal O}^{st}(F)}
f\int_{{E}_u(F)}\kappa_1(t_*)|\omega_{{E}_u}|
|\omega_{\overline{\Cal O}^{st}}|,
\endsplit
$$
where the measures are as in (0.6).
\endproclaim
As a consequence of this, it will follow that
spurious divisors do not contribute to the Igusa
expansion of (0.1) or (0.2).

The proof that $e(\theta,\beta)\leq 1$ in the case of
$G_2$ will now be sketched. From (0.1) we know that the
germs associated to the
$(\beta -1)$-regular unipotent orbits ${\Cal O}$ get
``matched'' with the terms in the Igusa expansion of
homogeniety degree $\beta -1$ in $\lambda$. Since both sides depend
on $X$ and $\lambda$ we might write such a comparison briefly
as
$
\Gamma_\beta (\lambda,X,f)=I_\beta(\lambda,X,f),
$
the left side denoting the sum of the normalized $\kappa$-stable
germs of homogeniety degree $\beta -1$ and the right side
denoting the corresponding term in the Igusa expansion
(0.2), (0.3). If we replace $X$ in both sides by $cX$ then both sides
should have the same degree of homogeniety in $c$ as well. However,
by computing explicitly the dependence of the term $I_\beta(\lambda,X,f)$
in the Igusa expansion on $X$ we will see, using the case-by-case
computation in \S 5, that $e(\theta,\beta)\geq 2$ implies that the two sides
have different degrees of homogeniety in $c$, a contradiction.

It remains then to analyze the integrals
over the fundamental divisors. We cannot compute these integrals
but, in most cases, we can determine their dependence on $\lambda$ 
and the tangent direction of $C$, and this is enough for our purposes.
On the right-hand side of
formula 0.6(a) we find what is essentially the factor
$$
\Delta_{0}(X)=\lim_{\lambda\rightarrow 0}
{{\Delta_{loc}(\gamma_{\lambda,G},\gamma_{\lambda,H})}\over
D_{G/H}(\gamma_{\lambda,G},\gamma_{\lambda,H})},
$$
in the notation of \cite{H3, \S 2}.
The local translation invariance property in the hypotheses to
Hales' reduction lemma 5.19 below holds for
$\kappa(t_*)\Delta_0(X)$ as we will see (this is carried out in \S\S6-8
below, which we compute a quantity which is essentially
$\Delta_0(X)$).
Therefore, we leave the $\kappa$-stable ${\Cal O}$-germ
unchanged if we replace $X$ by $X+Y$, where $Y$ is sufficiently small
and as in lemma 5.19.
We may now apply lemma 5.18 to obtain the matching for functions
whose support meets the stable orbit ${\Cal O}^{st}(F)$ but no
other unipotent orbit.
(This is where \cite{LS2} is implicitly used.)
This descent argument, remarkably, allows us to
avoid doing the difficult, more detailed calculations necessary for a 
direct determination of the dependence of 
$\Delta_{0}(X)\kappa(t_\sigma)$ on $X$.
\par
As we've already mentioned, the $0$-regular orbit has already been matched
\cite{S2}.
The $1$-regular $\kappa$-stable germ of $G_2$ matches with the 
$1$-regular stable germ of $PGL(3)$ on cubic Cartans and with the
$1$-regular stable germ of $(SL(2)\times SL(2))/\{\pm 1\}$ 
on biquadratic Cartans.
The stable germ associated to the 
trivial unipotent orbit of $(SL(2)\times SL(2))/\{\pm 1\}$ 
matches with the 2-regular $\kappa$-stable
germ of $G_2$ on $T$, where $T$ and $\kappa$ determine
$H_{T,\kappa}=(SL(2)\times SL(2))/\{\pm 1\}$. The stable germ associated to the 
trivial unipotent orbit of $PGL(3)$ matches with the 3-regular $\kappa$-stable
germ of $G_2$ on $T$, where $T$ and $\kappa$ determine $H_{T,\kappa}=PGL(3)$.
These matchings are in accordance with the fact that the $2-$regular
and $3-$regular unipotent orbits of $G_2(\overline{F})$ 
are endoscopically induced from the trivial orbits of $A_1\times A_1$
and $A_2$, respectively \cite{Spal, ch. III, table on page 221}.
Finally, the matching for the  trivial orbit follows from the work
of Rogawski, Howe, and Harish-Chandra \cite{Ro}.
\par
Though I've tried to make this paper self-contained (outside of
the usual background in group theory \cite{Bor} and 
algebraic geometry \cite{M}), familiarity
with the ideas in \cite{L}, \cite{LS3}, \cite{H1}, and \cite{H3} is assumed.
\vskip .1in 
Acknowledgement: I'd like to thank Magdy Assem, Dan Bump,
Caroline Grant, Roland Martin, and especially Tom Hales for many 
enlightening conversations on the subject of this paper.
\enddocument

